% !TeX spellcheck = en_GB
\ifcsname SlidesDistr\endcsname%
\documentclass[handout,aspectratio=169]{beamer}
\else%
\documentclass[aspectratio=169]{beamer}
\fi%
\usepackage{fontspec}
\usepackage[T1]{fontenc}
\usepackage{amsmath}
\usepackage{amsfonts}
\usepackage{amssymb}
\usepackage{graphicx}
\usepackage{csquotes}
\usepackage{booktabs}
\usepackage{multicol}
\usepackage{enumerate}
\usepackage{microtype}
\usepackage[labelfont=bf,font={small}]{caption}
\usepackage{hyperref}
\usepackage{booktabs}
\usepackage{subcaption}
\usepackage{fancyhdr}
\usepackage{pdfpages}
\usepackage{siunitx}
\usepackage{tikz}
\usepackage{mdframed}

\defaultfontfeatures{Mapping=tex-text}
\newfontfamily\symbolfont{Symbola}
\newfontfamily\quotefont{GentiumPlus}

\usepackage[sorting=none]{biblatex}
\addbibresource{../bibliography.bib}

\author{Chris Eliasmith}


\renewcommand{\vec}[1]{{\mathbf{#1}}}
\newcommand{\mat}[1]{{\mathbf{#1}}}
\newcommand{\T}{\ensuremath{\mathsf{T}}}
\renewcommand{\epsilon}{\varepsilon}
\renewcommand{\phi}{\varphi}

% Tango color palette
\definecolor{butter1}{HTML}{FCE94F}
\definecolor{butter2}{HTML}{EDD400}
\definecolor{butter3}{HTML}{C4A000}
\definecolor{orange1}{HTML}{FCAF3E}
\definecolor{orange2}{HTML}{F57900}
\definecolor{orange3}{HTML}{CE5C00}
\definecolor{chocolate1}{HTML}{E9B96E}
\definecolor{chocolate2}{HTML}{C17D11}
\definecolor{chocolate3}{HTML}{8F5902}
\definecolor{chameleon1}{HTML}{8AE234}
\definecolor{chameleon2}{HTML}{73D216}
\definecolor{chameleon3}{HTML}{4E9A06}
\definecolor{skyblue1}{HTML}{729FCF}
\definecolor{skyblue2}{HTML}{3465A4}
\definecolor{skyblue3}{HTML}{204A87}
\definecolor{plum1}{HTML}{AD7FA8}
\definecolor{plum2}{HTML}{75507B}
\definecolor{plum3}{HTML}{5C3566}
\definecolor{scarletred1}{HTML}{EF2929}
\definecolor{scarletred2}{HTML}{CC0000}
\definecolor{scarletred3}{HTML}{A40000}
\definecolor{aluminium1}{HTML}{EEEEEC}
\definecolor{aluminium2}{HTML}{D3D7CF}
\definecolor{aluminium3}{HTML}{BABDB6}
\definecolor{aluminium4}{HTML}{888A85}
\definecolor{aluminium5}{HTML}{555753}
\definecolor{aluminium6}{HTML}{2E3436}

\definecolor{violet}{HTML}{AA305C}
\definecolor{uwyellow}{HTML}{FDD433}
\definecolor{background}{HTML}{F9F9F6}
\definecolor{text}{HTML}{000000}

\definecolor{uweng1}{HTML}{D1B2EE}
\definecolor{uweng2}{HTML}{BF33DE}
\definecolor{uweng3}{HTML}{8001B3}
\definecolor{uweng4}{HTML}{56048A}

\setbeamercolor{title}{fg=violet}
\setbeamercolor{frametitle}{fg=black}
\setbeamercolor{structure}{fg=aluminium5}
\setbeamercolor{normal text}{fg=text}

\setbeamertemplate{navigation symbols}{}
\setbeamertemplate{footline}[frame number]

\hypersetup{%
	colorlinks=false,% hyperlinks will be black
	urlbordercolor=aluminium4,% hyperlink borders will be red
	pdfborderstyle={/S/U/W 0.5}% border style will be underline of width 1pt
}

\makeatletter
\newcommand{\superimpose}[2]{%
	{\ooalign{{#1}\hidewidth\cr{#2}\hidewidth\cr}}}
\makeatother
\newcommand{\SolidCircle}[2]{\superimpose{\color{#1}\symbolfont ⬤}{\textbf{\color{white}#2}}\hspace{1em}}
\newcommand{\OPlus}{\SolidCircle{chameleon3}{\kern0.75pt+}}
\newcommand{\OMeh}{\SolidCircle{uwyellow}{~}}
\newcommand{\OMinus}{\SolidCircle{scarletred3}{\kern2.25pt--}}

\newcommand{\hl}[1]{\colorbox{uwyellow}{{\color{black}\textbf{#1}}}}

\newcommand{\ImageSources}[1]{%
	\begin{columns}%
		\column{1.1\textwidth}%
		\raggedright%
		\tiny\color{aluminium4}%
		\setlength\lineskip{1em}%
		\textbf{Image Sources.}	{#1}%
	\end{columns}}

\newcommand{\ColorRect}[3]{{\color{#1}\rule{#2}{#3}}}
\setbeamertemplate{headline}{\ColorRect{black}{\textwidth}{4pt}\newline\ColorRect{uweng1}{0.25\textwidth}{4pt}\ColorRect{uweng2}{0.25\textwidth}{4pt}\ColorRect{uweng3}{0.25\textwidth}{4pt}\ColorRect{uweng4}{0.25\textwidth}{4pt}}

\newcommand{\MakeTitle}{%
	\vspace{0.5cm}%
	{\textbf{\inserttitle}}\\[0.5cm]%
	\insertauthor\\[0.5cm]%
	\insertdate\\%
	\begin{itemize}
		\item Slide design: Andreas Stöckel
		\item Content: Terry Stewart, Andreas Stöckel, Chris Eliasmith
	\end{itemize}
 	\includegraphics[width=7cm]{../assets/uwlogo_eng.pdf}%
}

\newcommand{\handwritingframe}{%
	\begin{frame}
		\begin{columns}
			\column{\paperwidth}
			\includegraphics{../assets/handwriting_lines.pdf}
		\end{columns}
	\end{frame}	
}

\newcommand{\imageframe}[1]{%
	\setbeamertemplate{navigation symbols}{}%
	\begin{frame}[plain,noframenumbering]%
		\begin{tikzpicture}[remember picture,overlay]%
		\node[at=(current page.center)] {%
			\includegraphics[width=\paperwidth]{#1}%
		};%
		\end{tikzpicture}%
	\end{frame}%
}

\newcommand{\videoframe}[3][mp4]{%
	\begin{frame}[plain,noframenumbering]%
		\hypersetup{%
			pdfborderstyle={/S/U/W 0}% border style will be underline of width 1pt
		}%
		\begin{tikzpicture}[remember picture,overlay]%
		\node[at=(current page.center)] {%
			\includegraphics[width=\paperwidth]{{video/#2_#3}.jpg}%
		};%
		\node[at=(current page.center)] {%
			\ifcsname SlidesDistr\endcsname%
				\href{https://youtu.be/#3}{\includegraphics[width=2cm]{../assets/play_button.pdf}}%
			\else%
				\href{video/#2_#3.#1}{\includegraphics[width=2cm]{../assets/play_button.pdf}}%
			\fi%
		};%
		\end{tikzpicture}%
	\end{frame}%
}

\newcommand{\includevideo}[4][mp4]{%
	\begingroup%
	\hypersetup{%
		pdfborderstyle={/S/U/W 0}% border style will be underline of width 1pt
	}%
	\begin{tikzpicture}%
	\node (A) {%
		\includegraphics[width=#4]{{video/#2_#3}.jpg}%
	};%
	\node[at=(A.center)] {%
		\ifcsname SlidesDistr\endcsname%
			\href{https://youtu.be/#3}{\includegraphics[width=2cm]{../assets/play_button.pdf}}%
		\else%
			\href{video/#2_#3.#1}{\includegraphics[width=2cm]{../assets/play_button.pdf}}%
		\fi%
	};%
	\end{tikzpicture}%
	\endgroup%
}

\newcommand{\backupbegin}{
	\newcounter{finalframe}
	\setcounter{finalframe}{\value{framenumber}}
	\setbeamertemplate{footline}{}
}

\newcommand{\backupend}{
	\setcounter{framenumber}{\value{finalframe}}
}


\date{September 27 \& Octoboer 2, 2023}
\title{SYDE 556/750 \\ Simulating Neurobiological Systems \\ Lecture 5: Feed-Forward Transformation}

\begin{document}
	
	\begin{frame}{}
		\vspace{0.5cm}
		\begin{columns}[c]
			\column{0.6\textwidth}
			\MakeTitle
			\column{0.4\textwidth}
			\includegraphics[width=\textwidth]{media/yellow_butterfly_by_albert_bierstadt_small.jpg}
		\end{columns}
	\end{frame}

	\begin{frame}{Introduction}
    \begin{columns}[c]
			\column{0.6\textwidth}
        \begin{itemize}
          \setlength\itemsep{0.5cm}
          \item We've only talked about representation til now
          \begin{itemize}
            \setlength\itemsep{0.25cm}
            \item What about computation?
          \end{itemize}
          \item We start by focusing on the state of a network after learning and development
          \item A kind of hypothesis testing and generation 
        \end{itemize}
      \column{0.4\textwidth}
        \includegraphics[width=\textwidth]{media/dalle_child_adult.png}
        \raggedleft\small DALL-E AI Generated Art, 2022
    \end{columns}
	\end{frame}

  \begin{frame}{NEF Principle 2: Transformation}
		\centering
		\includegraphics[scale=1.175]{media/transformation_fully_connected.pdf}
		\begin{mdframed}
			\textbf{NEF Principle 2 -- Transformation}\\
			Connections between populations describe \emph{transformations} of neural representations. Transformations are functions of the variables represented by neural populations.
		\end{mdframed}
	\end{frame}
	
	\begin{frame}{A Tale of Two Populations (I)}
		\hspace*{-1.2cm}%
		\includegraphics<1>[width=1.175\textwidth]{media/transformation_01.pdf}%
		\includegraphics<2>[width=1.175\textwidth]{media/transformation_02.pdf}%
	\end{frame}

	\begin{frame}{Communication Channel Experiment: Same input signal}
		\includegraphics[width=\textwidth]{media/two_populations_decoding.pdf}
	\end{frame}

	\begin{frame}{A Tale of Two Populations (II)}
		\hspace*{-1.2cm}%
		\includegraphics<1>[width=1.175\textwidth]{media/transformation_03.pdf}%
		\includegraphics<2>[width=1.175\textwidth]{media/transformation_04.pdf}%
	\end{frame}

	\begin{frame}{Communication Channel Experiment: Populations in series}
		\includegraphics<1>[width=\textwidth]{media/two_populations_decoding.pdf}
		\includegraphics<2>[width=\textwidth]{media/two_populations_decoding_series.pdf}
	\end{frame}

	\begin{frame}{Computing Synaptic Weights: Step 1 -- Encoding Matrix}
		\hspace*{-1.2cm}%
		\includegraphics<1>[width=1.175\textwidth]{media/transformation_04.pdf}%
		\includegraphics<2>[width=1.175\textwidth]{media/transformation_05.pdf}%
	\end{frame}

	\begin{frame}{Computing Synaptic Weights: Step 2 -- Scaled Encoding Matrix}
		\hspace*{-1.2cm}%
		\includegraphics<1>[width=1.175\textwidth]{media/transformation_06.pdf}%
	\end{frame}

	\begin{frame}{Computing Synaptic Weights: Step 3 -- $\mat W = \mat E \mat D$}
		\hspace*{-1.2cm}%
		\includegraphics<1>[width=1.175\textwidth]{media/transformation_07.pdf}%
	\end{frame}

  \begin{frame}{Computational Complexity}
    \begin{itemize}
      \setlength\itemsep{0.25cm}
        \item Weights - multiplying $\vec a \in \mathbb{R}^{n}$ with $\mat W \in \mathbb{R}^{m \times n}$ is $\mathcal{O}(n m)$ i.e., $\approx\mathcal{O}(n^2)$
        \item Decoding - $\hat{\vec x} = \mat D \vec a$ is $\mathcal{O}(d n)$
        \item Encoding - $\vec J = \mat E \hat{\vec x} + \vec J_\mathrm{bias}$ is $\mathcal{O}(d m)$
        \item Encoding/Decoding - $\mathcal{O}(d (n + m))$ or $\approx\mathcal{O}(d n)$ for $n = m$
        \item So if $d$ is small we get a linear complexity $\mathcal{O}(n)$
        \item Therefore, sequential decoding and re-encoding saves a lot of time compared to using actual synaptic weights
        \item One reason why Nengo is so fast compared to other SNN simulators
      \end{itemize}
  \end{frame}

	\begin{frame}{Computing Functions}
		\hspace*{-1.2cm}%
		\includegraphics<1>[width=1.175\textwidth]{media/transformation_02.pdf}%
		\includegraphics<2>[width=1.175\textwidth]{media/transformation_08.pdf}%
		\includegraphics<3>[width=1.175\textwidth]{media/transformation_09.pdf}
		\begin{overlayarea}{\textwidth}{1cm}
			\only<3->{\hl{Function Decoder} $
				\mat D^f = \left( (\mat A \mat A^\T + N \sigma^2 \mat I)^{-1} \mat A \mat Y^\T \right)^\T \,, \quad
				\text{where } \big ( \mat Y \big )_{ik} = \big( f(\vec x_k) \big)_i 
			$}
		\end{overlayarea}
	\end{frame}

	\begin{frame}{Decoding Functions -- Using a Few Neurons}
		\includegraphics<1>[width=\textwidth]{media/function_decodings_10.pdf}%
	\end{frame}

	\begin{frame}{Decoding Functions -- Using More Neurons}
		\includegraphics<1>[width=\textwidth]{media/function_decodings_20.pdf}%
		\includegraphics<2>[width=\textwidth]{media/function_decodings_50.pdf}%
		\includegraphics<3>[width=\textwidth]{media/function_decodings_100.pdf}%
	\end{frame}
	
	\begin{frame}{Computing Functions -- Weight Matrix}
		\hspace*{-1.2cm}%
		\includegraphics<1>[width=1.175\textwidth]{media/transformation_10.pdf}
		\begin{overlayarea}{\textwidth}{1cm}
			\centering
			$\mat W^f = \mat E \mat D^f$
		\end{overlayarea}
	\end{frame}

  \begin{frame}{Recipe for any feedforward transformation}
    \begin{enumerate}
      \setlength\itemsep{0.25cm}
        \item Define encoding for two populations (input/output)
        \item Write the transformation with the represented input variables
        \item Solve for the $\mat D^f$ for that transformation for the input population
        \item Sub 3. into the encoding for the output variable
      \end{enumerate}
  \end{frame}

	\begin{frame}{Computing Multivariate Functions}
		\centering
		\includegraphics{media/network_legend_3.pdf}\\[0.5cm]
		\begin{columns}[t]
			\column{0.33\textwidth}
			\centering
			\hl{Linear Superposition}\\[0.25cm]
			$W^{f_1} \vec a_1(\vec x) + W^{f_2} \vec a_2 (\vec y)$
			\includegraphics{media/network_a.pdf}
			\pause
			\column{0.33\textwidth}
			\centering
			\hl{Nonlinear Functions}\\[0.25cm]
			Multi-dimensional $\vec z$
			\includegraphics{media/network_c.pdf}
			\pause
			\column{0.33\textwidth}
			\centering
			\textbf{(Dendritic Computation)}\\[0.25cm]
			Exploit dendritic nonlinearity
			\includegraphics{media/network_d.pdf}
		\end{columns}
	\end{frame}

	\begin{frame}{Computing Multivariate Functions -- Linear Superposition}
		\begin{columns}[c]
			\column{0.33\textwidth}
			\centering
			\hl{Linear Superposition}\\[0.25cm]
			\includegraphics{media/network_a.pdf}
			\column{0.66\textwidth}
			\includegraphics[width=\textwidth]{media/multivariate_addition.pdf}
		\end{columns}
 	\end{frame}


	\begin{frame}{Computing Multivariate Functions -- Multiplication}
		\begin{columns}[c]
			\column{0.33\textwidth}
			\centering
			\hl{Nonlinear Functions}\\[0.25cm]
			Multi-dimensional $\vec z$
			\includegraphics{media/network_c.pdf}\\[0.25cm]
			\begin{overlayarea}{\textwidth}{2.75cm}
				\only<2->{Multiplication is useful...
				\begin{itemize}
					\item Gating of signals
					\item Attention effects
					\item Binding
					\item Statistical inference				
				\end{itemize}}
			\end{overlayarea}
			\column{0.66\textwidth}
			\includegraphics[width=\textwidth]{media/multivariate_multiplication.pdf}
		\end{columns}
	\end{frame}


	\backupbegin
	
	\begin{frame}[noframenumbering]{Image sources}
		\small
		\textbf{Title slide}\\\enquote{Yellow Butterfly}\\Author: Albert Bierstadt, circa 1890.\\From \href{https://commons.wikimedia.org/wiki/File:Yellow_Butterfly_by_Albert_Bierstadt.jpeg}{Wikimedia}.
	\end{frame}


	\backupend
	
\end{document}
